\documentclass{article}
\usepackage{graphicx}
\usepackage{mathtools}
\usepackage{amssymb}

\title{IMC 2026 P1}
\author{W.F Mathis BILLA}
\date{November 2026}

\begin{document}

\maketitle

\textbf{P1:} Show that the equation $\cos(\cos x)=\sin(\sin x)$ has no real solution.



Let's call the equation $(A)$: $\cos(\cos x)=\sin(\sin x)$.

Suppose $x \in \mathbb{R}$ is a solution of $(A)$.

Using the identity $\cos X=\sin\left(\frac{\pi}{2}-X\right)$ we can write:
$$
\sin\left(\frac{\pi}{2}-\cos x\right)=\sin(\sin x)
$$
and since $\sin X=\sin Y \iff X=Y+2k\pi$ or $X=\pi-Y+2k\pi$ $(k \in \mathbb{Z})$, we get
$$
\cos x+\sin x=\frac{\pi}{2}-2k\pi \quad (1)
\qquad \textbf{or} \qquad
\sin x-\cos x=\frac{\pi}{2}+2k\pi \quad (2)
\qquad (k \in \mathbb{Z}).
$$



\textbf{(1).}
\begin{align*}
(1) &\iff \sqrt{2}\cdot\sin\left(x+\frac{\pi}{4}\right)=\frac{\pi}{2}-2k\pi \\
    &\iff \sin\left(x+\frac{\pi}{4}\right)=\frac{\sqrt{2}}{2}\left(\frac{\pi}{2}-2k\pi\right)
\end{align*}
We know $-1\le\sin X\le 1$ $\;\forall X\in\mathbb{R}$. So we have
\begin{align*}
-1 \le \frac{\sqrt{2}}{2}\left(\frac{\pi}{2}-2k\pi\right) \le 1
&\iff -\sqrt{2}\le \frac{\pi}{2}-2k\pi\le\sqrt{2} \\
&\iff \frac{1}{2}\left(\frac{\pi}{2}-\sqrt{2}\right)\le k\pi\le \frac{1}{2}\left(\frac{\pi}{2}+\sqrt{2}\right) \\
&\iff \frac{\pi-2\sqrt{2}}{4\pi}\le k\le\frac{\pi+2\sqrt{2}}{4\pi}.
\end{align*}
Since $\pi>3$ and $2\sqrt{2}<3$ (because $8<9$), we have $\pi-2\sqrt{2}>0$. Also $\pi+2\sqrt{2}<4\pi$, because $2\sqrt{2}<3\pi$. So
$$
0<\frac{\pi-2\sqrt{2}}{4\pi}\le k\le\frac{\pi+2\sqrt{2}}{4\pi}<1,
$$
and no integer $k$ satisfies this.


\textbf{(2).}
Since $\sin\left(x-\frac{\pi}{4}\right)=\frac{\sqrt{2}}{2}(\sin x-\cos x)$,
\begin{align*}
(2) &\iff \sin\left(x-\frac{\pi}{4}\right)=\frac{\sqrt{2}}{2}\left(\frac{\pi}{2}+2k\pi\right).
\end{align*}
As before, $-1\le \frac{\sqrt{2}}{2}\left(\frac{\pi}{2}+2k\pi\right)\le 1$, so
\begin{align*}
-\sqrt{2}\le \frac{\pi}{2}+2k\pi\le\sqrt{2}
&\iff \frac{-\pi-2\sqrt{2}}{4\pi}\le k\le\frac{-\pi+2\sqrt{2}}{4\pi}.
\end{align*}
Since $2\sqrt{2}<3<\pi$ and $\pi+2\sqrt{2}<4\pi$, we get
$$
-1<\frac{-\pi-2\sqrt{2}}{4\pi}\le k\le\frac{-\pi+2\sqrt{2}}{4\pi}<0,
$$
and again no integer $k$ satisfies this.



 Hence neither $(1)$ nor $(2)$ can hold for any $k\in\mathbb{Z}$, which contradicts the existence of $x$. We conclude that $(A)$ has no solution.

\\
\\


\end{document}